% Sections 7 and 8: limitations and conclusion.
% Claims and status: a synthesis of Sections 2-5; no new numbers. The served comparisons are
% those of Section 4.5 (committed; single sessions except the frontier, the composed stack and the
% served certified head, cert-served).
\section{Limitations}\label{sec:discussion}\label{sec:limits}

\paragraph{The error model is an assumption} Agreement with \Rstock{} rests on an assumed model of a closed-source kernel's accumulation; neither of the two models covers BF16 split-$K$ partial sums (Section~\ref{sec:headcontracts}). The crafted-input probes behind the \modelhopper{} used products with one exact operand and reached neither subnormals nor other GPUs, so neither model is established for all inputs. The replayed positions cannot tell the two models apart (Section~\ref{sec:rstock-share}). The engine's BF16-output kernel hides its accumulator, so the error measured on its FP32-output variant is indirect evidence \ev{kernel-error}. The int8 kernel is checked per compiled variant, not proved to compute the modelled arithmetic (Section~\ref{sec:kernel-design}).

\paragraph{Replay and serving} Recomputation and fallback figures come from plain-decoding replay at batches of up to 16 positions, then consecutive ones, and the engine check from 64 prompts under the \modelcons{} only. The served runs, on one greedy workload, did not time the in-graph fallbacks, so attributing the shortfall to them is a hypothesis.

\paragraph{Scope} Transport was refuted for two drafters and the bound families and tilings of Table~\ref{tab:transport}, in a replay that favours it by treating draft logits as exact. A drafter trained to reproduce the target's head input is untested. The measurements concern one model, engine commit and GPU, with text-only prompts. The contract and proofs carry over to any head that rounds an FP32 accumulation to BF16 and breaks ties by index; the error model and the fallback shares must be re-established for each kernel and GPU.

\section{Conclusion}\label{sec:conclusion}
A verifier does not need every logit, but the token it needs is the engine's own, rounding and tie rule included, not the exact-arithmetic winner. On this engine an int8 copy of the head finds it with one or two exact recomputations per decision, under a named bound on the kernel's error; served, it gains 1.2--4.6\% throughput at one request with identical tokens and loses on DFlash above one request and on MTP at 64, where fallbacks become common, and 0.6--1.2\% gated off. Bounds carried over from the drafters' heads fail in every variant tested: their head inputs are far from the target's. The design that makes the certified head safe, a fast path held to the engine's output with its kernel as fallback, is not specific to the head. On this hybrid model the head is a small target (6\% of a plain step at batch 128, the state kernel 42\%; Appendix~\ref{app:profile}); other changes gain more against stock (three sessions): the fold with FA4 attention adds 8.4--9.3\% to DFlash's per-request rate at $c=1$--4 and 12.4--13.7\% to throughput at 8--32, and delayed, 16-capped prefill batches give MTP 1.18--1.32 times the best non-speculative arm's throughput at 48--128, at a later first token.
